\chapter{General Introduction}

\section{Chapter Introduction}
This opening chapter establishes the research problem addressed in the thesis and defines the academic scope of the work. It introduces the internship-discovery context, identifies the limitations of purely lexical and LLM-dependent approaches, states the research questions, and summarizes the main contributions. The chapter concludes by explaining how the remaining chapters are organized.

\section{Context and Motivation}
Over the last decade, internship discovery has been transformed by the digital shift in the global labor market \cite{faliagka2012online}. Traditional channels, such as newspaper classifieds and physical noticeboards, have largely been replaced by a broad ecosystem of online platforms. Students now navigate global platforms such as LinkedIn and Indeed, specialized university career portals, and niche industry boards. This shift has expanded access to a large volume of listings \cite{kurekova2015using}. However, this abundance also creates a significant data-management challenge. Internship listings are often semi-structured, inconsistent, and expressed through platform-specific conventions. Two platforms, or even two employers, may describe the same role using different terms. A ``Junior Java Developer'' and a ``Software Engineering Intern'' might require an identical skill set, yet a traditional keyword search treats them as separate entities. In addition, the data is often noisy, with boilerplate text, inconsistent formatting, and non-standard labels \cite{tamburri2020data}. For recruitment platforms and universities, organizing this continuous flow of unstructured data from automated web scraping, manual JSON uploads, and other sources is a substantial task. This challenge motivates the need for a lightweight approach that can transform noisy listings into a structured semantic knowledge base, enabling precise discovery without the overhead of manual curation or heavyweight computational infrastructure \cite{boselli2018classifying}.

\section{Problem Statement}
This thesis addresses two related problems in the ingestion and retrieval of internship data: the semantic gap and structural heterogeneity. Existing job-posting systems commonly follow one of two approaches, each of which creates limitations for the target use case:


\begin{itemize}[leftmargin=1.4em,itemsep=0.25em,topsep=0.25em]

\item Rigid Lexical Matching: Many platforms rely on exact keyword matching, such as BM25. This approach is limited in its ability to handle synonyms, technical jargon, and relationships between concepts. As a result, users may miss relevant listings because their search terms do not precisely match the employer's vocabulary \cite{baeza2011modern}.

\item High-Cost LLM Dependency: The opposite extreme is to delegate the problem entirely to a Large Language Model (LLM). While LLMs are effective at interpreting meaning, relying on them to process thousands of scraped documents in real time is often financially expensive and can introduce substantial latency \cite{menghani2023efficient}. The ingestion process itself is similarly fragmented. Data from partner institutions often arrives in JSON formats that lack any fixed schema, forcing teams into a repetitive cycle of manual mapping. This creates the need for a unified system that can automatically clean, structure, and semantically enrich data from diverse sources while remaining efficient enough to run on ordinary commodity hardware \cite{giabelli2021skills}.

\end{itemize}
\section{Research Questions}
The study is guided by four research questions:


\begin{itemize}[leftmargin=1.4em,itemsep=0.25em,topsep=0.25em]

\item RQ1: How can a tiered pipeline architecture balance extraction accuracy and computational cost for internship listings?

\item RQ2: How well can lightweight transformer models such as DistilBERT map non-standard job titles onto a standardized professional taxonomy such as ESCO?

\item RQ3: How much does a hybrid retrieval strategy---one that combines lexical (BM25) and semantic (HNSW) search---improve the relevance and recall of internship discovery compared to traditional keyword search alone?

\item RQ4: Which heuristic-based mapping strategies can automate the ingestion of inconsistently formatted JSON data into a unified schema?

\end{itemize}
\section{Contributions}
This thesis makes four primary contributions:


\begin{itemize}[leftmargin=1.4em,itemsep=0.25em,topsep=0.25em]

\item A Tiered Semantic Enrichment Pipeline: The thesis introduces a modular architecture that combines three layers: deterministic, rule-based extraction (Layer A), lightweight neural mapping (Layer B), and a resource-budgeted LLM fallback for the most complex cases (Layer C).

\item A Hybrid Search Infrastructure: The thesis implements a retrieval system using an in-memory HNSW index and BM25 to support hybrid lexical and semantic search.

\item Automated Ingestion Heuristics: The thesis develops a flexible JSON processing module that uses fuzzy key matching and content-based inference to reduce manual schema mapping for multi-format user uploads.

\item Empirical Evaluation: The thesis presents an evaluation using a real-world dataset of 327 listings and reports the observed trade-offs between semantic enrichment, retrieval quality, and computational cost.

\end{itemize}

\section{Thesis Structure}
The thesis is organized into seven chapters. Chapter 1 introduces the research context, problem statement, research questions, contributions, and overall structure. Chapter 2 presents the background, related work, and theoretical foundations of the study, including information extraction, occupational taxonomies, text representation, semantic retrieval, hybrid search, and lightweight model strategies. Chapter 3 describes the technological environment and development tools that support the implementation, including the client-side application layer, API layer, database layer, transformer-based semantic tools, and retrieval infrastructure. Chapter 4 presents the methodology, covering the system overview, semantic schema, label space, and evaluation protocol. Chapter 5 describes the system design and implementation, including data collection, preprocessing, enrichment, storage, indexing, caching, and API design. Chapter 6 presents the reported results, adds a structured procedure for completing HR@k and NDCG@k analysis, and preserves the comparison section as a completion scaffold rather than as an additional empirical claim. Chapter 7 combines discussion, limitations, conclusion, and future-work material, with the ethical, legal, and privacy section retained as a completion scaffold because it was listed in the original thesis structure. The appendices preserve the original appendix headings and identify the material still to be completed.

\section{Chapter Conclusion}
This chapter has defined the motivation, research problem, questions, and contributions of the thesis. It has also clarified the boundaries of the study: the proposed system is evaluated as an internship aggregation and retrieval framework for internship listings, while incomplete external-platform comparison evidence is treated as work that must be completed with measured results.
